\documentclass{amsart}
% For dealing with references we use the comment environment
\usepackage{verbatim}
\newenvironment{reference}{\comment}{\endcomment}
%\newenvironment{reference}{}{}
\newenvironment{slogan}{\comment}{\endcomment}
\newenvironment{history}{\comment}{\endcomment}

% For commutative diagrams we use Xy-pic
\usepackage[all]{xy}

% We use 2cell for 2-commutative diagrams.
\xyoption{2cell}
\UseAllTwocells

% We use multicol for the list of chapters between chapters
\usepackage{multicol}

% This is generally recommended for better output
\usepackage{lmodern}
\usepackage[T1]{fontenc}

% For cross-file-references

% Package for hypertext links:
\usepackage{hyperref}

% For any local file, say "hello.tex" you want to link to please

% Theorem environments.
%
\theoremstyle{plain}
\newtheorem{theorem}[subsection]{Theorem}
\newtheorem{proposition}[subsection]{Proposition}
\newtheorem{lemma}[subsection]{Lemma}

\theoremstyle{definition}
\newtheorem{definition}[subsection]{Definition}
\newtheorem{example}[subsection]{Example}
\newtheorem{exercise}[subsection]{Exercise}
\newtheorem{situation}[subsection]{Situation}

\theoremstyle{remark}
\newtheorem{remark}[subsection]{Remark}
\newtheorem{remarks}[subsection]{Remarks}

\numberwithin{equation}{subsection}

% Macros
%
\def\lim{\mathop{\mathrm{lim}}\nolimits}
\def\colim{\mathop{\mathrm{colim}}\nolimits}
\def\Spec{\mathop{\mathrm{Spec}}}
\def\Hom{\mathop{\mathrm{Hom}}\nolimits}
\def\Ext{\mathop{\mathrm{Ext}}\nolimits}
\def\SheafHom{\mathop{\mathcal{H}\!\mathit{om}}\nolimits}
\def\SheafExt{\mathop{\mathcal{E}\!\mathit{xt}}\nolimits}
\def\Sch{\mathit{Sch}}
\def\Mor{\mathop{\mathrm{Mor}}\nolimits}
\def\Ob{\mathop{\mathrm{Ob}}\nolimits}
\def\Sh{\mathop{\mathit{Sh}}\nolimits}
\def\NL{\mathop{N\!L}\nolimits}
\def\CH{\mathop{\mathrm{CH}}\nolimits}
\def\proetale{{pro\text{-}\acute{e}tale}}
\def\etale{{\acute{e}tale}}
\def\QCoh{\mathit{QCoh}}
\def\Ker{\mathop{\mathrm{Ker}}}
\def\Im{\mathop{\mathrm{Im}}}
\def\Coker{\mathop{\mathrm{Coker}}}
\def\Coim{\mathop{\mathrm{Coim}}}

% Boxtimes
%
\DeclareMathSymbol{\boxtimes}{\mathbin}{AMSa}{"02}

%
% Macros for moduli stacks/spaces
%
\def\QCohstack{\mathcal{QC}\!\mathit{oh}}
\def\Cohstack{\mathcal{C}\!\mathit{oh}}
\def\Spacesstack{\mathcal{S}\!\mathit{paces}}
\def\Quotfunctor{\mathrm{Quot}}
\def\Hilbfunctor{\mathrm{Hilb}}
\def\Curvesstack{\mathcal{C}\!\mathit{urves}}
\def\Polarizedstack{\mathcal{P}\!\mathit{olarized}}
\def\Complexesstack{\mathcal{C}\!\mathit{omplexes}}
% \Pic is the operator that assigns to X its picard group, usage \Pic(X)
% \Picardstack_{X/B} denotes the Picard stack of X over B
% \Picardfunctor_{X/B} denotes the Picard functor of X over B
\def\Pic{\mathop{\mathrm{Pic}}\nolimits}
\def\Picardstack{\mathcal{P}\!\mathit{ic}}
\def\Picardfunctor{\mathrm{Pic}}
\def\Deformationcategory{\mathcal{D}\!\mathit{ef}}

\setlength{\emergencystretch}{3em}
\begin{document}
\title[Stacks Project, Tag 02OA]{373 --- Ideal-power cohomology and ideal powers of cohomology are pro-isomorphic under properness}
\maketitle
\noindent Copyright (C) 2005--2025 Johan de Jong. Stacks Project authors. Source: \href{https://stacks.math.columbia.edu/tag/02OA}{Tag 02OA}.

\noindent Editorial difficulty note (not author text). Difficulty H2: The proof must control kernels and transition images uniformly, not merely compare completed limits. Finite graded Rees-module generation and generation of its kernel module yield shifted transition bounds, allowing explicit reverse maps and verification of both pro-system compositions..

\noindent Editorial provenance note (not author text). History: excerpt from revision \texttt{a0444\allowbreak 6e57e\allowbreak c1fbc\allowbreak 252a8\allowbreak 71afc\allowbreak ec775\allowbreak 2fb28\allowbreak 07b14}, coherent.tex, lines 5423--5538. Reference presentation was adapted; original-author context and core support blocks were retained or added. The mathematical wording is unchanged. Licensed under GNU FDL 1.2 or later; no invariant sections or cover texts. See ../licenses/COPYING. Context blocks reproduce author definitions and supporting results; they are not additional counted problems.

\begin{lemma}
\label{lemma-cohomology-powers-ideal-application}
Let $A$ be a Noetherian ring.
Let $I \subset A$ be an ideal.
Let $f : X \to \Spec(A)$ be a proper morphism.
Let $\mathcal{F}$ be a coherent sheaf on $X$.
Then for every $p \geq 0$ there exists an integer $c \geq 0$
such that
\begin{enumerate}
\item the multiplication map
$I^{n - c} \otimes H^p(X, I^c\mathcal{F}) \to H^p(X, I^n\mathcal{F})$
is surjective for all $n \geq c$,
\item the image of $H^p(X, I^{n + m}\mathcal{F}) \to H^p(X, I^n\mathcal{F})$
is contained in the submodule $I^{m - e} H^p(X, I^n\mathcal{F})$
where $e = \max(0, c - n)$ for $n + m \geq c$, $n, m \geq 0$,
\item we have
$$
\Ker(H^p(X, I^n\mathcal{F}) \to H^p(X, \mathcal{F})) =
\Ker(H^p(X, I^n\mathcal{F}) \to H^p(X, I^{n - c}\mathcal{F}))
$$
for $n \geq c$,
\item there are maps $I^nH^p(X, \mathcal{F}) \to H^p(X, I^{n - c}\mathcal{F})$
for $n \geq c$ such that the compositions
$$
H^p(X, I^n\mathcal{F}) \to 
I^{n - c}H^p(X, \mathcal{F}) \to
H^p(X, I^{n - 2c}\mathcal{F})
$$
and
$$
I^nH^p(X, \mathcal{F}) \to
H^p(X, I^{n - c}\mathcal{F}) \to
I^{n - 2c}H^p(X, \mathcal{F})
$$
for $n \geq 2c$ are the canonical ones, and
\item the inverse systems $(H^p(X, I^n\mathcal{F}))$ and
$(I^nH^p(X, \mathcal{F}))$ are pro-isomorphic.
\end{enumerate}
\end{lemma}

\begin{proof}
Write $M_n = H^p(X, I^n\mathcal{F})$ for $n \geq 1$ and
$M_0 = H^p(X, \mathcal{F})$ so that we have maps
$\ldots \to M_3 \to M_2 \to M_1 \to M_0$. Setting
$B = \bigoplus_{n \geq 0} I^n$, then $M = \bigoplus_{n \geq 0} M_n$
is a finite graded $B$-module, see
Lemma \href{https://stacks.math.columbia.edu/tag/02O8}{lemma cohomology powers ideal times F (Tag 02O8)}.
Observe that the products
$B_n \otimes M_m \to M_{m + n}$, $a \otimes m \mapsto a \cdot m$
are compatible with the maps in our inverse system in the sense
that the diagrams
$$
\xymatrix{
B_n \otimes_A M_m \ar[r] \ar[d] & M_{n + m} \ar[d] \\
B_n \otimes_A M_{m'} \ar[r] & M_{n + m'}
}
$$
commute for $n, m' \geq 0$ and $m \geq m'$.

\medskip\noindent
Proof of (1). Choose $d_1, \ldots, d_t \geq 0$ and $x_i \in M_{d_i}$
such that $M$ is generated by $x_1, \ldots, x_t$ over $B$.
For any $c \geq \max\{d_i\}$ we conclude that
$B_{n - c} \cdot M_c = M_n$ for $n \geq c$ and we conclude (1) is true.

\medskip\noindent
Proof of (2). Let $c$ be as in the proof of (1). Let $n + m \geq c$.
We have $M_{n + m} = B_{n + m - c} \cdot M_c$.
If $c > n$ then we use $M_c \to M_n$ and the compatibility of products
with transition maps pointed out above to conclude that
the image of $M_{n + m} \to M_n$ is contained in $I^{n + m - c}M_n$.
If $c \leq n$, then we write $M_{n + m} = B_m \cdot B_{n - c} \cdot M_c =
B_m \cdot M_n$ to see that the image is contained in $I^m M_n$.
This proves (2).

\medskip\noindent
Let $K_n \subset M_n$ be the kernel of the map $M_n \to M_0$. The
compatibility of products with transition maps pointed out above
shows that $K = \bigoplus K_n \subset M$ is a graded $B$-submodule.
As $B$ is Noetherian and $M$ is a finitely generated graded $B$-module,
this shows that $K$ is a finitely generated graded $B$-module.
Choose $d'_1, \ldots, d'_{t'} \geq 0$ and $y_i \in K_{d'_i}$
such that $K$ is generated by $y_1, \ldots, y_{t'}$ over $B$.
Set $c = \max(d'_i, d'_j)$. Since $y_i \in \Ker(M_{d'_i} \to M_0)$
we see that $B_n \cdot y_i \subset \Ker(M_{n + d'_i} \to M_n)$.
In this way we see that $K_n = \Ker(M_n \to M_{n - c})$ for $n \geq c$.
This proves (3).

\medskip\noindent
Consider the following commutative solid diagram
$$
\xymatrix{
I^n \otimes_A M_0 \ar[r] \ar[d] &
I^nM_0 \ar[r] \ar@{..>}[d] &
M_0 \ar[d] \\
M_n \ar[r] &
M_{n - c} \ar[r] &
M_0
}
$$
Since the kernel of the surjective arrow $I^n \otimes_A M_0 \to I^nM_0$
maps into $K_n$ by the above we obtain the dotted arrow and the
composition $I^nM_0 \to M_{n - c} \to M_0$ is the canonical map.
Then clearly the composition $I^nM_0 \to M_{n - c} \to I^{n - 2c}M_0$
is the canonical map for $n \geq 2c$. Consider the composition
$M_n \to I^{n - c}M_0 \to M_{n - 2c}$. The first map
sends an element of the form $a \cdot m$
with $a \in I^{n - c}$ and $m \in M_c$
to $a m'$ where $m'$ is the image of $m$ in $M_0$.
Then the second map sends this to $a \cdot m'$ in $M_{n - 2c}$ and
we see (4) is true.

\medskip\noindent
Part (5) is an immediate consequence of (4) and the definition of
morphisms of pro-objects.
\end{proof}
\end{document}
